

Diese Diplomarbeit befasst sich mit der Entwicklung der webbasierten Lernplattform \textbf{LearnHub}, die als Partnerprojekt für die HTL Leonding konzipiert wurde. Dabei standen die methodische Trennung von Lern- und Testphasen, sowie die Implementierung effizienter Lernstrategien im Vordergrund.

Im Rahmen der Projektentwicklung konnten wir uns besonders im technischen Bereich neues Know-how aneignen. Durch die Arbeit mit modernen Frameworks im Fullstack-Bereich sowie die Anbindung einer zentralen REST-Schnittstelle konnten wir unser bereits bestehendes Wissen vertiefen. Im Laufe des Entwicklungsprozesses haben wir diese Fertigkeiten kontinuierlich ausgebaut. Auch die technische Umsetzung der Lernstrategien, wie die Kombination aus dem Leitner-System und Spaced Repetition, erforderte einiges an Planung, um eine automatisierte Wiederholungslogik zu ermöglichen. Rückblickend sind wir sehr zufrieden mit dem Endergebnis von LearnHub, da der Übungsmodus durch direktes Feedback und der Prüfungsmodus durch eine realitätsnahe Simulation überzeugen.

Trotzdem gibt es Aspekte, die wir im Nachhinein anders angegangen wären. Besonders bei den Usability-Tests besteht Verbesserungspotenzial. Aufgrund von Zeitmangel konnten diese nicht im geplanten Ausmaß durchgeführt werden. Eine frühzeitige Einbindung von Nutzertests hätte geholfen, den Userflow noch flüssiger zu gestalten und die Bedienbarkeit der Plattform weiter zu optimieren.

In der Zukunft bietet LearnHub ein funktionales Fundament für weitere Erweiterungen.